\section[Adoption]{A realistic starting point}
\sectionframe{09}{A realistic starting point}
  {What should I change in my repository first?}
  {\faSeedling}

% ---------------------------------------------------- other skill packages
% \pkgcard{k}{owner}{repository}{head colour}{author}{what it is}{its steps}{its rule for tests}: card k
% (0 to 3), 33.6 mm wide; the green box at its foot holds the rule, level with the other cards' rules
\newcommand\pkgcard[8]{%
  \pgfmathsetmacro\pkx{#1*36.8}%
  \begin{scope}
    \clip[rounded corners=1.2mm] (\pkx,0) rectangle ++(33.6,-45.4);
    \fill[white] (\pkx,0) rectangle ++(33.6,-45.4);
    \fill[#4] (\pkx,0) rectangle ++(33.6,-9.2);
    \fill[greentint] (\pkx,-32.6) rectangle ++(33.6,-12.8);
    \fill[testgreen] (\pkx,-32.6) rectangle ++(1.2,-12.8);
  \end{scope}
  \draw[udsgreylight,line width=0.5pt,rounded corners=1.2mm] (\pkx,0) rectangle ++(33.6,-45.4);
  \node[anchor=north west,inner sep=0pt,font=\ttfamily\fontsize{7}{8.4}\selectfont,text=white!70!#4,
    align=left] at (\pkx+1.8,-1.5) {#2/\\\textcolor{white}{\bfseries #3}};
  \node[pkg who] at (\pkx+1.8,-10.8) {#5};
  \node[pkg what] at (\pkx+1.8,-14.8) {#6};
  \node[pkg text] at (\pkx+1.8,-22.6) {#7};
  \node[pkg text,text width=28.2mm] at (\pkx+3.6,-34.4) {\textcolor{testgreen}{\faCheck}\enspace #8};}
\begin{frame}[source={GitHub, 3 Nov 2026: obra/superpowers v7.0.1 · multica-ai/andrej-karpathy-skills · open-gsd/gsd-core v1.20.0 · mattpocock/skills v1.3.1}]{Other skill packages, one shared core}
  \begin{tikzpicture}[x=1mm,y=1mm,
      pkg who/.style={anchor=north west,inner sep=0pt,font=\scriptsize,text=udsgrey,text width=30.4mm,align=left},
      pkg what/.style={anchor=north west,inner sep=0pt,font=\footnotesize\bfseries,text=udsdark,
        text width=30.4mm,align=left,execute at begin node={\hyphenpenalty=10000\raggedright}},
      pkg text/.style={anchor=north west,inner sep=0pt,font=\scriptsize,text=udsdark,text width=30.4mm,
        align=left,execute at begin node={\hyphenpenalty=10000\raggedright}}]
    \useasboundingbox (0,0) rectangle (144,-54);
    \pkgcard{0}{mattpocock}{skills}{udsdark}{Matt Pocock}{small skills you combine}
      {grill, spec, tickets, implement, review}{\fn{/tdd}: red before green, at agreed seams}
    \pkgcard{1}{obra}{superpowers}{udsblue}{Jesse Vincent}{a whole method in 15 skills}
      {brainstorm, plan, execute, review, finish}{no production code without a failing test first}
    \pkgcard{2}{multica-ai}{andrej-karpathy-skills}{udsblue}{from Karpathy's notes}{one \fn{CLAUDE.md}, four rules}
      {think first, keep it simple, change only what is asked, set checkable goals}
      {a bug fix starts with a test that reproduces the bug}
    \pkgcard{3}{open-gsd}{gsd-core}{udsblue}{continues Get Shit Done}{a whole project in phases, 72 commands}
      {discuss, plan, execute, verify, ship}{tests mapped to requirements before any code; TDD optional}
    % what all four share
    \fill[greentint] (0,-48) rectangle (144,-54);
    \node[anchor=west,inner sep=0pt,font=\footnotesize\bfseries,text=testgreen] at (2,-51)
      {\faCheckDouble\enspace in all four, the check is set before the code is written};
  \end{tikzpicture}
  \note{\textbf{41:50--42:50}\par
    Pocock's collection is one of several; three others are worth a look.\par
    Superpowers, by Jesse Vincent, is a whole method, from brainstorming to finishing the
      branch; its test-driven skill allows no production code without a failing test.\par
    The Karpathy skills are one CLAUDE.md with four rules, written by Jiayuan Zhang from Andrej
      Karpathy's notes on how models go wrong in code; the last rule turns a task into a goal
      that can be checked, such as a test that reproduces a bug.\par
    GSD Core, the continuation of Get Shit Done, runs a whole project in spec-driven phases; it
      maps test coverage to the requirements before any code, and full TDD is an option.\par
    They differ in how much of the process they take over; Pocock's README prefers small skills
      you combine to systems that own the process.\par
    What they share is section 7: the check is decided before the code.\par
    As with any skill, read them before installing and note the version.}
\end{frame}

% --------------------------------------------------------------- checklist
\begin{frame}{Checklist for your repository}
  \setbeamertemplate{itemize item}{\raisebox{-0.1ex}{\textcolor{udsblue}{\faSquare[regular]}}}
  \setlength\leftmargini{5mm}
  \begin{columns}[T,onlytextwidth]
    \begin{column}{0.3\textwidth}
      \textbf{\color{udsdark}Groundwork}\par\vskip1mm
      \begin{itemize}\footnotesize\setlength\itemsep{2mm}
        \item a git repository; commit before each agent task
        \item the environment in one file, with a lockfile
        \item Ruff and mypy configured, run in CI
      \end{itemize}
    \end{column}
    \begin{column}{0.31\textwidth}
      \textbf{\color{udsblue}Context}\par\vskip1mm
      \begin{itemize}\footnotesize\setlength\itemsep{2mm}
        \item \fn{AGENTS.md}: commands, boundaries, pointers
        \item \fn{GLOSSARY.md}: the key terms
        \item \fn{CODING\_STANDARDS.md}: what review checks
        \item \fn{docs/adr/}: decisions people reopen
      \end{itemize}
    \end{column}
    \begin{column}{0.33\textwidth}
      \textbf{\color{udsred}Control}\par\vskip1mm
      \begin{itemize}\footnotesize\setlength\itemsep{2mm}
        \item one fast check on synthetic data
        \item tests agreed before the code
        \item raw data read-only; confidential data out of reach
        \item sandbox, network and approvals set in advance
        \item a named reviewer for the code and the science
      \end{itemize}
    \end{column}
  \end{columns}
  \note{\textbf{42:50--43:40}\par
    Groundwork: a git repository with a commit before every agent task, an environment described
      in one file with a lockfile, and Ruff and mypy running in CI.\par
    Context: AGENTS.md, the glossary, the coding standards and ADRs, so that requests can stay
      short.\par
    Control: one reliable check, tests agreed before the code, data and secrets out of reach,
      deliberate sandbox and approval settings, and a person who reviews both the code and the
      science.\par
    None of this needs a new tool, and most of it helps human collaborators too.}
\end{frame}

% -------------------------------------------------------------- discussion
\begin{frame}[noframenumbering]{Questions and discussion}
  % one full-width column, so the list starts where the columns of other slides do
  \begin{columns}[T,onlytextwidth]
    \begin{column}{\textwidth}
      \begin{itemize}\setlength\itemsep{2.4mm}
        \item Which of your projects is not under git yet?
        \item Which repository would you try this on first, and who reviews it?
        \item Which decision in your project deserves an ADR?
        \item Which single command tells you that your code still works?
        \item Which data must never be within an agent's reach?
      \end{itemize}
    \end{column}
  \end{columns}
  \note{\textbf{43:40--44:30}\par
    Close with concrete questions: projects still without version control, a first repository to try
      and its reviewer, a decision worth recording, the one command that shows the code still
      works, and data that must stay out of reach.\par
    Templates and references follow in the appendix.}
\end{frame}
